% Summary for the mistral-rse profile: Research Software Engineer, Mistral AI
% (Paris). Written against that posting's first responsibility -- "accelerate
% researchers by owning the complex parts of large-scale pipelines and
% delivering robust internal tooling" -- and says nothing the experience and
% education sections below do not back up.
%
% Deliberately absent: Kubernetes and SLURM. The posting asks for them and the
% claim would not be true; the reproducible-container and GPU-node work is what
% is offered against that requirement instead. Do not "improve" this by adding
% an orchestrator.
%
% Three verbs, not one. The middle of this paragraph once read "I own the CUDA
% kernels ... , the institute's CI/CD, its DGX A100/H200 nodes and its
% vision-language models" -- one verb stretched over four unlike things, and it
% fitted none of them well: you write kernels, you run nodes, you serve models.
% Write/run/serve says what the work actually is at no extra length.
%
% The paragraph is about Fraunhofer IPA, deliberately. It used to spend its
% whole second half on personal repositories -- "my renderers, my ML code and
% this page's pipeline are public" -- which is roughly half the space a reader
% gives this CV, handed to spare-time work. The open-source claim did not have
% to go with them: the fixes to the institute's toolchain are *the job*, not a
% hobby, so saying it that way keeps the claim, moves it onto paid work, and
% buys the line the vision-language clause needed. github.com/Kataglyphis in
% the header, the Hobbies column and the footer's typeset-with credit carry
% what is left of the personal side; they do not need a sentence here too.
%
% "not into a local patch" is load-bearing and has been lost to a length trim
% once already. The claim is a habit with a consequence, not "passionate about
% open source" -- the contrast is what makes it a practice rather than a taste.
% Backed by the Experiences bullet on upstreamed fixes, which no longer names
% the two PRs it used to, on request. Worth its line for this employer in
% particular: Mistral ships open-weight models, and the Fraunhofer bullets are
% careful to say open-weight where that is what they mean.
%
% Every item here is named, not explained: each one is a bullet in
% section_experience immediately below, and a summary that re-explains the page
% under it is the first thing a reader skips. The one clause that goes further
% -- "that replaced its hand-labeling" -- is there because it is an outcome,
% not a restatement, and because it is the only thing in the paragraph that
% says what "faster" bought anyone.
%
% One page is a hard requirement and this paragraph is the first thing to break
% it: three lines fit, four push the Languages/Hobbies/Publications row onto
% page two -- which has already happened here twice. The ceiling is about 345
% rendered characters in either language, and German reaches it sooner, which
% is why it says "die das Labeling übernehmen" where the English can afford
% "that replaced its hand-labeling". Rebuild BOTH languages after any edit.
% No GitHub URL here: the header already carries it.
\IfLanguageName{english}{
    \par{I make other people's research run faster: at Fraunhofer IPA I write the CUDA kernels and Python bindings behind the team's methods, run its CI/CD and DGX A100/H200 nodes, and serve the vision-language models that replaced its hand-labeling. Where the tooling we depend on is open source, my fixes go upstream, not into a local patch.}
}{
    \par{Ich mache die Forschung anderer schneller: Am Fraunhofer IPA schreibe ich CUDA-Kernels und Python-Bindings hinter den Verfahren des Teams, betreibe CI/CD und DGX-A100/H200-Knoten und hoste die Vision-Language-Modelle, die das Labeling übernehmen. Wo unsere Werkzeuge Open Source sind, gehen meine Fixes upstream statt in lokale Patches.}
}
